\newpage
\section*{Questão 2 — Análise Politrópica e Consumo Específico}
\label{sec:q_2}
\vspace{{10pt}}
\begin{{minipage}}{{\textwidth}}
\small
\paragraph{{Enunciado:}}
\begin{{verbatim}}
Um motor a gasolina é alimentado com um combustível cujo poder calorífico é de 42.000 kJ/kg. 
As pressões no interior do cilindro, quando o pistão está a 5\% e a 75\% do curso de compressão, 
são de 1,2 bar e 4,8 bar, respectivamente.

Admitindo que o processo de compressão obedece à lei politrópica: $PV^{1,3} = constante$ 
e desprezando os efeitos de biela-manivela considerando $V$ variando linearmente com o curso, determine:
\begin{enumerate}[label=\alph*)]
\item A razão de compressão do motor;
\item Se a relação da eficiência indicada real do motor com a eficiência do ciclo ar-padrão ($\eta_i/\eta_{ciclo}$), 
é de 60\%, calcule o consumo específico indicado de combustível em g/kWh.
\end{enumerate}

\end{{verbatim}}
\end{{minipage}}
\vspace{{15pt}}
\subsection*{Solução}
\noindent\textbf{Dados:}\\[5pt]
\begin{center}
\begin{tabular}{|c|r|c|} \hline
\textbf{Parâmetro} & \textbf{Valor} & \textbf{Unidade} \\ \hline\hline
$PCI$ & 42000 & kJ/kg \\ \hline
$P_a$ & 1.2 & bar \\ \hline
$P_b$ & 4.8 & bar \\ \hline
$y_a$ & 0.05 & -- \\ \hline
$y_b$ & 0.75 & -- \\ \hline
$n_{poly}$ & 1.3 & -- \\ \hline
$\eta_{rel}$ & 0.60 & -- \\ \hline
$k$ & 1.4 & -- \\ \hline
\end{tabular}
\end{center}\\[10pt]

\paragraph{Resolução:}

\subsubsection*{Razão de Volumes Intermediários ($R_v$)}
\noindent A partir da lei politrópica da compressão $P_A V_A^n = P_B V_B^n$ [cite: 13], determinamos a razão volumétrica exata entre as posições de 5\% e 75\% do curso do pistão[cite: 12].
\[ R_{v} = \left(\frac{P_{b}}{P_{a}}\right)^{\frac{1}{n_{poly}}} \approx 2.90 \]
\vspace{6pt}

\subsubsection*{Razão de compressão do motor ($r$)}
\noindent A posição do pistão $y$ relaciona-se com o volume através da função linear $V(y) = V_c + (1-y)V_d$[cite: 13]. Substituindo os volumes $V_a$ e $V_b$ na razão $R_v$ e sabendo que $r = 1 + V_d/V_c$, isolamos a razão de compressão global do motor.
\[ r = \frac{R_{v} - 1}{1 - y_a - R_{v} \left(1 - y_{b}\right)} + 1 \approx 9.51 \]
\vspace{6pt}

\subsubsection*{Eficiência do ciclo ar-padrão (Otto)}
\noindent Como se trata de um motor a gasolina , a eficiência teórica é modelada pelo ciclo Otto ideal (ar-padrão)[cite: 14], que depende exclusivamente da razão de compressão calculada no passo anterior.
\[ \eta_{ciclo} = 1 - r^{1 - k} \approx 0.5938 \text{ ou } 59.38\% \]
\vspace{6pt}

\subsubsection*{Eficiência indicada real ($\eta_i$)}
\noindent A eficiência indicada real representa a fração efetiva da energia do combustível convertida em trabalho sobre o pistão. O problema fornece que essa eficiência corresponde a 60\% da eficiência ideal[cite: 14].
\[ \eta_{i} = \eta_{ciclo} \cdot \eta_{rel} \approx 0.3563 \text{ ou } 35.63\% \]
\vspace{6pt}

\subsubsection*{Consumo Específico Indicado ($CEI$)}
\noindent O consumo específico quantifica a massa de combustível necessária para gerar 1 kWh de energia indicada[cite: 14]. Para realizar a conversão dimensional correta (de kJ para kWh e de kg para gramas), aplicamos o fator de conversão de $3.6 \times 10^6$.
\[ CEI = \frac{3.6 \times 10^6}{PCI \cdot \eta_{i}} \approx 240.56 \text{ g/kWh} \]
\vspace{6pt}

\newpage